\scanpage{60}
\mspage{29}
\margin{\deleted{\illegible{1}}multiplicités modulaires \inserted{$\widehat M_G$} de dim $2$}
\[
(74)\quad\left\{\begin{array}{ll}
\widehat M^!_{0,5}&\Gamma\text{ opère via }\Gamma\longrightarrow\mathfrak S_{\widehat I_{\alpha_0}}\ (\simeq\mathfrak S_5)\\
\widehat M^!_{1,2}&\Gamma\text{ --- }\quad\Gamma\longrightarrow\mathfrak S_{\widehat I_{\alpha_0}}\simeq\mathbb Z/2\mathbb Z\\
\widehat M^!_{0,4}\times\widehat M^!_{0,4}&\Gamma\text{ --- }\quad\Gamma\longrightarrow\deleted{\unsure{\prod_{\alpha\in\{\alpha_0,\alpha_1\}}}}\mathfrak S_{\deleted{\unsure{\widehat I}}\bigl(\widehat I_{\alpha_0}\amalg\widehat I_{\alpha_1}}\simeq\mathfrak S_8\\
\widehat M^!_{0,4}\times\widehat M^!_{1,1}&\Gamma\text{ --- }\quad\Gamma\longrightarrow\mathfrak S^{\deleted{\unsure{4|*}}}_{\widehat I_{\alpha_0}}\simeq\mathfrak S_4\\
\widehat M^!_{1,1}\times\widehat M^!_{1,1}&\Gamma\text{ opère trivialement.}
\end{array}\right.
\]

\uline{NB} \deleted{\unsure{à}} Dans le cas $3^\circ$) de $\widehat M_{0,4}\times\widehat M_{0,4}$, $\Gamma$ \deleted{opère}\inserted{s'envoie} dans $\mathfrak S_8$ par l'intermédiaire de
\[
\begin{array}{rcl}
\Gamma&\longrightarrow&\mathbb Z/2\mathbin{\cdot}(\mathfrak S_4\times\mathfrak S_4)\\
&&\text{prod. }\frac12\text{ direct}
\end{array}
\]

On peut considérer que par rapport aux schémas modulaires \inserted{$\widehat M_G$} de dim $1$, \deleted{\unsure{qui sont}}\inserted{ils sont deux :}
\margin{(cas élémentaires!)}
\[
(75)\quad\left\{\begin{array}{ll}
\widehat M^!_{0,4}&\Gamma\text{ opère via }\Gamma\longrightarrow\mathfrak S_{I_{\alpha_0}}\simeq\mathfrak S_4\\
\widehat M^!_{1,1}&\Gamma\text{ opère trivialt},
\end{array}\right.
\]
les ``cas \unsure{''|``}nouveaux'' de (74) sont $\widehat M^!_{0,5}$ et $\widehat M^!_{1,2}$.

Cela fait en tous essentiellement quatre cas ``élémentaires'' à étudier à fond, pour le cas \unsure{des|de} dim. modulaires $\leq2$. J'ai l'impression qu'une compréhension plus ou moins complète de la structure des $\Pi_1$ \inserted{discrets (i.e. dans le contexte analytique)} impliqués dans ces cas, \unsure{donnerait|donnera} la clef pour une compréhension complète dans le cas général.
\margin{59}
